\chapter{Goldbach deserts are not prime gaps: the height game}
\label{ch:gaps}

\begin{chapterabstract}
A prime gap of length $Q$ just below $\N$ forces $\g(\N)>Q$, so the
height game could in principle free-ride on the prime-gap records. This
chapter shows it should not: blocking only the \emph{prime} offsets is
strictly cheaper than blocking an interval, and the separation is now
certified, not heuristic. Two constructions --- $\g=1\,157\,341$ at
$2\,480$ digits and $\g=1\,134\,871$ at $2\,692$ digits, from
independent pipelines --- exceed the longest prime gap ever certified
with proven endpoints ($1\,113\,106$, between primes of $18\,662$
digits) using integers roughly seven times shorter than that gap's own
endpoints. The second construction additionally carries a fully
certified \emph{local} prime gap of just $8\,970$: its Goldbach desert
is $126.5$ times longer than the prime-free interval it sits in, so the
desert owes essentially nothing to an ordinary gap. At these heights the
search is nearly free ($\Ee\approx6.6$) and the binding constraint
becomes the primality certificate for the partition --- making the
height record an index of primality-proving technology rather than of
number-theoretic difficulty --- and we close with the one construction
idea that would remove that wall too.
\end{chapterabstract}

\section{The transfer, and its one-way sign}\label{sec:transfer5}

Proposition~\ref{prop:transfer} sends prime gaps to Goldbach deserts:
if $[\N-Q,\N-2]$ is prime-free then $\g(\N)>Q$ (given any partition at
all). The converse direction is where the money is, and it fails
constructively: a desert needs $\N-q$ composite only for the
$\pi(Q)$ \emph{primes} $q<Q$, a set of density $1/\ln Q$ inside the
interval the gap must clear entirely. Three certified exhibits measure
the separation.

\paragraph{Exhibit 1: a desert that is a gap in disguise.} The
primorial $m_{43}=13\,082\,761\,331\,670\,030$ of
Remark~\ref{rem:concrete} has $\g(m_{43})=89$, and the nearest prime
below it is $m_{43}-89$ itself: this small desert coincides with an
ordinary prime-free interval. At toy scale the two phenomena are
entangled.

\paragraph{Exhibit 2: a desert 57 times its gap.} The first-generation
$237$-digit record ($\g=109\,621$) already separates them: the
committed verifier proves the interval $[\N-1926,\,\N-2]$ prime-free
while $\N-1927$ is a probable prime ($1927$ is composite, so it is no
summand). Ordinary gap information therefore certifies only
$\g>1\,926$; the actual value is $57\times$ larger, entirely on the
back of composite offsets that the desert may ignore.

\paragraph{Exhibit 3: the certified ratio.} The $2\,692$-digit
construction below carries \emph{both} numbers as certificates: desert
$1\,134\,871$, local prime gap $8\,970$. Ratio: $126.5$.

\section{Benchmarks from the gap world}\label{sec:benchmarks}

The longest prime gap known with \emph{proven} endpoints has length
$1\,113\,106$ and sits between primes of $18\,662$ digits; the longest
known with probable-prime endpoints has length
$16\,045\,848$~\cite{GapList}. Two framing facts. First, certifying a
gap is brutally expensive at scale: the endpoints need primality
proofs, and every interior integer a compositeness witness ---
which is why the proven-endpoint record lags the PRP record by an order
of magnitude. Second, gap \emph{merit} (length divided by $\ln$ of the
endpoint) measures how anomalous a gap is: the largest merit ever
observed is $\approx42$, and Cram\'er--Granville-type
heuristics~\cite{GranvilleCramer} cap plausible merits near a few
hundred. Hiding a million-length desert inside a genuine prime gap at
$200$ digits would need merit $\approx2\,180$ --- fantasy. Deserts at
practical sizes \emph{must} come from covering primes, not intervals;
this chapter's constructions are that statement in certified form.

The challenge the height game set itself: beat $1\,113\,106$ ---
a desert longer than any interval ever proven prime-free --- with
integers small enough to certify on a lab box. (The challenge benchmark
was set at $1\,113\,137$, marginally above the record gap; both
constructions clear both numbers.)

\section{The 2\,480-digit record: $\g=1\,157\,341$}\label{sec:g1}

The CPU pipeline's height construction uses $759$ odd congruence
classes plus parity ($M\approx10^{2474.6}$), leaving $|U|=2\,844$
residual primes below the target $Q=1\,113\,137$, and scans the
progression $\N=\N_0+tM$. At $2\,480$ digits the exponent is tiny ---
$p_1|U|\approx7.7$, expected $\approx$$2.2\cdot10^3$ progression
elements, $\approx1\,800$ \emph{tested} elements under selective
testing --- and the hit arrived at $t=14\,544$, the $443$rd tested
candidate. Certification is where the effort lives:

\begin{itemize}
\item every one of the $89\,840$ primes $p<\g$ proven a non-summand:
$86\,714$ by covering divisor, $3\,126$ by failed strong
probable-prime test, all re-derived by the standalone verifier;
\item $\g=1\,157\,341$ prime by trial division; the complement
$\N-\g$ (a $2\,480$-digit integer) proven prime by a $245$-step ECPP
certificate ($\approx$1.7 hours to generate), re-verified from scratch
by the independent pure-Python checker ($7\,036$ s), \emph{plus}
APR-CL as an algorithmically distinct second proof ($7\,409$ s);
\item an independent PARI full scan of all $89\,841$ primes
$p\le\g$: zero violations ($27.4$ minutes).
\end{itemize}

Result: $\g(\N)=1\,157\,341$, exceeding the certified-gap record by
$4\%$ with an integer $7.5\times$ fewer digits than that gap's
endpoints. The desert record and the gap record are now formally
decoupled: ours needed no interval at all.

\section{The 2\,692-digit mega-desert, and its certified local gap}\label{sec:megagap}

The GPU-session pipeline independently built a second height
construction --- $813$ odd congruences plus parity, $2\,692$ digits,
found at $k=14$ with measured $\hat\Ee=6.56$ against predicted $6.6$
--- with
\[
\g(\N)=1\,134\,871,
\]
proven by $88\,239$ offset witnesses (parity $1$, congruence divisor
$85\,476$, trial divisor $895$, strong base-2 $1\,867$) and an
ECPP-certified complement, the $253$-step certificate re-verified by
the independent checker with its binding to $\N-\g$ asserted.

What makes this construction scientifically distinctive is the
\emph{extra} pair of certificates: the nearest primes around $\N$ are
pinned at $\N-7\,329$ and $\N+1\,641$ --- each carrying its own
validated ECPP certificate, with every interior integer deterministically
composite --- so the ordinary prime gap containing $\N$ has length
exactly $8\,970$, fully certified. The desert-to-gap ratio,
\[
\frac{\g(\N)}{\text{local prime gap}}
=\frac{1\,134\,871}{8\,970}=126.5,
\]
is the thesis's cleanest single number for the claim that Goldbach
deserts are not prime gaps: $99.2\%$ of this desert's length lies
\emph{outside} any prime-free interval, bought by covering prime
offsets alone. In digit terms the separation reads: a $2\,692$-digit
integer carries a $1.13$-million desert, while an interval of that
length free of primes is only known (unproven!) around integers of
$\sim$$18\,000$ digits --- the sparse-cover advantage is about
$7\times$ in digits today and widens with the target.

\section{The economics of height: certification is the wall}\label{sec:certwall}

The cost model explains why the height game is played at thousands of
digits. $\Ee=|U|\cdot\mathrm{boost}/\ln\N$ \emph{falls} as $\N$ grows
--- doubling the digits roughly halves the exponent at fixed cover
strength --- so unlike the digit-frugal bounded games, the height game
optimises at whatever size makes the search trivial:
$\Ee\approx6.6$ at $2\,700$ digits means hits arrive by the dozen. What
does not get cheaper is proving the partition:

\begin{itemize}
\item ECPP generation and validation are minutes at $200$ digits,
$\approx$2 hours (plus 2 hours of independent re-verification, plus
2 hours of APR-CL) at $2\,480$--$2\,692$ digits, and scale roughly
quartically in the digit count;
\item at $\sim$4\,000 digits, days; far beyond that, the certificate
--- not the desert --- is the record's binding constraint with present
tools.
\end{itemize}

The height record is therefore an index of primality-proving
technology. Extrapolating with the cover economics: present ECPP
practice supports $\g\approx2.5$ million (the next committed rung);
distributed fastECPP-class certification ($\sim$$50\,000$-digit
proofs) would support $\g\approx5\cdot10^7$; a certification
breakthrough to $10^6$-digit general-form proofs would carry
$\g$ toward $10^9$. The desert machinery itself is already there: even
the absolute-gap milestone --- beating the $16\,045\,848$ PRP-endpoint
gap, i.e.\ covering $\pi(1.6\cdot10^7)\approx1.03$ million prime
offsets with a $\sim$$20\,000$-digit modulus --- is a throughput
problem ($\sim$$0.2$ s per modular exponentiation at that width; a
distributed campaign $10^2$--$10^3\times$ one lab box), not a
conceptual one. The certificate is the only wall.

\section{Removing the wall: pick the provable prime first}\label{sec:twosided}

Every construction in this thesis searches over $\N$ and pays ECPP for
whichever complement $\N-\g$ falls out --- a random $2\,700$-digit
prime with no exploitable structure. Invert the design: \emph{fix the
complement first}, inside a family whose primality is cheap to
certify, and search the other side.

Concretely, take Proth-form candidates $P=j\cdot2^n+1$ with
$j<2^n$: Proth's theorem~\cite{Proth} certifies such a $P$ prime with a
\emph{single} modular exponentiation once a quadratic non-residue
witness is found --- minutes at $10^5$ digits, where general-form ECPP
is hopeless. For fixed $n$, the cover congruences
$\N\equiv a_r\pmod r$ with $\N=P+s$ ($s$ the intended least summand)
translate to congruences on $j$, since $2^n\bmod r$ is computable; the
search over $j$ is again a CRT progression, now on the \emph{certifiable}
side. A hit hands over the partition $\N=s+P$ with $P$'s certificate
essentially free, deleting the height game's only wall
(Pocklington-style variants~\cite{Pocklington} offer the same trade
with partially factored $P-1$). The open questions are the density of
Proth primes inside the cover-constrained $j$-progression and whether
the $n$-side and $s$-side constraints can be tuned jointly; nobody has
run the construction, and we flag it as the single most valuable open
trick in this problem family --- it would let the $\g$-ladder run to
Proth-search scale ($10^5$--$10^6$ digits, $\g\sim10^8$ and beyond)
with today's tools. Chapter~\ref{ch:conclusions} files it among the
open problems with its neighbours.
